\section{Semi-model structure and the homotopy hypothesis}
\label{sec:consequences}

We now apply the expansion theorem to the original category
\(\mathcal A=\Mod(\mathcal C)\), without a coskeletal cutoff.
The first consequence is the canonical structure of
\cite[Theorem~5.3.5]{Henry2016}.  We give its transfer argument explicitly
to distinguish cellular objects from their cofibrant retracts.

\subsection{The canonical semi-model structure}

Recall
\[
 I=\{I_j:\partial D_j\longrightarrow D_j\mid j\ge0\},
 \qquad
 J=\{J_j:D_j\longrightarrow D_{j+1}\mid j\ge0\}.
\]
For a set of maps \(K\), write \(K\text{-}\mathrm{cof}\)
for the maps with the left lifting property
against \(K\text{-}\mathrm{inj}\).

We use the left semi-model convention in which
cofibration--trivial-fibration factorization and lifting hold for all maps.
Trivial-cofibration--fibration factorization and lifting are required for maps
with cofibrant domain.  The proof below verifies these restrictions
explicitly.

\begin{corollary}\label{cor:semimodel}
For every Grothendieck coherator \(\mathcal C\), the category
\(\mathcal A\) admits a cofibrantly generated left semi-model structure
with
\[
 \begin{aligned}
  \text{weak equivalences}&=\mathcal W,\\
  \text{cofibrations}&=I\text{-}\mathrm{cof},\\
  \text{fibrations}&=J\text{-}\mathrm{inj},\\
  \text{trivial fibrations}&=I\text{-}\mathrm{inj}.
 \end{aligned}
\]
Here \(\mathcal W\) is the class defined by components and based
homotopy groups.  Cofibrant objects are exactly retracts of cellular
objects, and every object is fibrant.
\end{corollary}

\begin{proof}
Each \(J_j\) is a composite of two boundary attachments: an
\(I_j\)-attachment adds the parallel \(j\)-cell, and an
\(I_{j+1}\)-attachment adds the cell joining the two.  Consequently
\(J\subset I\text{-}\mathrm{cof}\).  The lifting calculation of
Lemma~\ref{lem:lifting-test} gives
\begin{equation}\label{eq:consequences-trivial-fibrations}
 I\text{-}\mathrm{inj}
   =\mathcal W\cap J\text{-}\mathrm{inj}.
\end{equation}
Indeed an \(I\)-injective map satisfies the exact-boundary test and
lifts \(J\); conversely, for a weak equivalence lifting \(J\), first
lift a disk up to one higher homotopy and then lift that homotopy to
correct the disk exactly.

The small object argument for \(I\) supplies a weak factorization
system
\[
 \bigl(I\text{-}\mathrm{cof},\ I\text{-}\mathrm{inj}\bigr)
\]
on all maps.  In particular, factoring \(\varnothing\to X\) shows
that every cofibrant \(X\) is a retract of a cellular object: the
cofibrancy of \(X\) provides a section of the resulting trivial
fibration onto \(X\).  The converse follows from closure of the left
lifting class under retracts.

Theorem~\ref{thm:main} extends from cellular starts to cofibrant starts.
To see this explicitly, choose a retract diagram
\(X\xrightarrow{u}Z\xrightarrow{v}X\), with \(vu=\id_X\) and
\(Z\) cellular.  Given \(a:D_j\to X\), form the two expansions
\[
 X^+=X\amalg_{D_j,a}D_{j+1},
 \qquad
 Z^+=Z\amalg_{D_j,ua}D_{j+1}.
\]
The maps \(u,v\), together with the identity on the new
\(D_{j+1}\), induce maps \(X^+\to Z^+\to X^+\) whose composite
is the identity.  Thus \(X\to X^+\) is a retract of \(Z\to Z^+\)
in the arrow category and belongs to \(\mathcal W\).

Every relative \(J\)-cell map with cofibrant domain is therefore a
weak equivalence.  Cofibrancy is preserved at each stage because every
\(J_j\) is an \(I\)-cofibration.  At successor stages use the
preceding paragraph, and at limit stages use the filtered-colimit
closure of \(\mathcal W\) from Lemma~\ref{lem:weak-closure}.
A set of simultaneous attachments can equally be attached in any
well-order, so this also covers the coproducts used by the small object
argument.  The small object argument for \(J\) consequently factors
every map with cofibrant domain as
\[
 X\xrightarrow{j}Z\xrightarrow{p}Y,
 \qquad
 j\in I\text{-}\mathrm{cof}\cap\mathcal W,
 \quad p\in J\text{-}\mathrm{inj}.
\]
As a relative \(J\)-cell map, \(j\) has the left lifting property
against every fibration.

Conversely, let \(i:X\to Y\) be an \(I\)-cofibration and a weak
equivalence, with \(X\) cofibrant.  Apply this factorization to \(i\),
writing \(i=pj\).  Two-out-of-three puts \(p\) in \(\mathcal W\),
so Equation~\eqref{eq:consequences-trivial-fibrations} makes \(p\)
\(I\)-injective.  A lift of \(i\) against \(p\) exhibits \(i\)
as a retract of \(j\).  Hence every trivial cofibration with
cofibrant domain lifts against every fibration.

These are precisely the required factorizations and lifting axioms:
cofibration--trivial-fibration factorization and lifting hold without
restriction, while trivial-cofibration--fibration factorization and
lifting are required for cofibrant domains.  Completeness and
cocompleteness follow from the algebraic description of \(\mathcal A\);
the retract and two-out-of-three axioms were proved in
Lemma~\ref{lem:weak-closure}.  Finally, the unit of a \(j\)-cell
extends it along \(J_j\).  Thus every map \(X\to 1\) is
\(J\)-injective, and every object is fibrant.
\end{proof}

\subsection{The comparison with spaces}

We use the cylinder-category comparison established in
\cite[Sections~5.1 and~5.3]{Henry2016}.  The relevant implication is
all-dimensional and has exactly the expansion assertion as its
additional hypothesis.  We spell out its application to the fixed
coherator.

The comparison first packages finite cellular models and their
cylinders into a small category.  Its completion recovers all models
of the same coherator.  Henry's recognition theorem then identifies
that completion with spaces up to Quillen equivalence.  The proof
below records the hypotheses at each of these three interfaces.

A cylinder category has cofibrations and weak equivalences satisfying
\cite[Definition~2.1.1]{Henry2016}; in particular, it has cofibration
pushouts, its weak equivalences satisfy two-out-of-six, every object
has a cylinder, and every trivial cofibration retracts.  Its completion
is the full category of set-valued presheaves taking the initial object
to a singleton and cofibration pushouts to pullbacks
\cite[Definition~2.4.1]{Henry2016}.

Here \(\Spaces\) denotes the usual model category of topological spaces,
with weak homotopy equivalences.

\begin{corollary}\label{cor:homotopy-hypothesis}
For every Grothendieck coherator \(\mathcal C\), the canonical
semi-model category \(\Mod(\mathcal C)\) is Quillen equivalent to
spaces.  More precisely, there is a geometric-realization/fundamental-
infinity-groupoid adjunction
\[
 |-|:\Mod(\mathcal C)\ \rightleftarrows\ \Spaces:\Pi_\infty
\]
which is a Quillen equivalence.  In particular,
\[
 \Mod(\mathcal C)[\mathcal W^{-1}]
       \simeq \operatorname{Ho}(\Spaces).
\]
Thus the Grothendieck homotopy hypothesis holds for these groupoids.
\end{corollary}

\begin{proof}
Let \(\mathcal A_{\mathrm{fin}}\) be a small skeleton of the full
subcategory of finitely cellular \(\mathcal C\)-groupoids, with finite boundary-cell
extensions as cofibrations and the restriction of \(\mathcal W\)
as weak equivalences.  Henry's comparison uses the following three
facts.

First, under the expansion assertion, \(\mathcal A_{\mathrm{fin}}\)
is a cylinder category
\cite[Proposition~5.3.6]{Henry2016}.  The finite nature of its
cylinders is explicit: a relative cylinder for the generating
cofibration \(\partial D_j\to D_j\) is
\[
 D_j\amalg_{\partial D_j}D_j
       \longrightarrow D_{j+1}\longrightarrow D_j.
\]
The first map is \(I_{j+1}\), and the source endpoint is \(J_j\);
the last map is induced by a unit.  The source endpoint is a weak
equivalence by Theorem~\ref{thm:main}, and hence so is the last map.
Finite relative cylinders are obtained by the finite cell induction
of that proposition.  Trivial cofibrations in this category have
retractions by Corollary~\ref{cor:semimodel}, since their domains are
cofibrant and all objects are fibrant.

Second, the free construction of the original coherator identifies
the completion of \(\mathcal A_{\mathrm{fin}}\) with
\(\Mod(\mathcal C)\), equipped with precisely the structure in
Corollary~\ref{cor:semimodel}
\cite[Proposition~5.3.7]{Henry2016}.  In this construction one begins
with finite globular sets and adjoins the operations in the given
free order of \(\mathcal C\).  Thus the identification retains the
chosen theory and its actual algebraic models.  The generating
cofibrations in the completion are the \(I_j\), and its generating
trivial cofibrations are the relative-cylinder endpoints \(J_j\).

Third, the marked object \(D_0\) makes
\(\mathcal A_{\mathrm{fin}}\) a cylinder coherator
\cite[Corollary~5.3.12]{Henry2016}.  In Henry's terminology this
means that the map from the free pre-cylinder category \(F_*\) on one object,
selecting \(D_0\), is a trivial cofibration and that the target is a
cylinder category \cite[Definition~5.1.1]{Henry2016}.
The hypotheses used in this implication are the
free cellular construction of \(\mathcal C\), the expansion
assertion, and the homotopical terminality of \(D_0\).  The first is
part of Definition~\ref{def:coherator}, the second is
Theorem~\ref{thm:main}, and the third follows here also from
Proposition~\ref{prop:shape-contractible}: the map \(D_0\to1\)
is \(I\)-injective by the boundary-filling characterization of
contractibility.  Every map \(A\to D_0\) therefore extends across
every finite cellular cofibration \(A\to B\); this gives
homotopical terminality inside \(\mathcal A_{\mathrm{fin}}\).
More explicitly, the projection
\(p:\mathcal A_{\mathrm{fin}}^{D_0}\to\mathcal A_{\mathrm{fin}}\)
in Henry's construction is an acyclic fibration by
\cite[Proposition~4.4.4]{Henry2016}.\footnote{The phrase
``trivial cofibration'' for this projection in the proof of
\cite[Lemma~5.3.11]{Henry2016} is a typographical error;
Proposition~4.4.4 gives an acyclic fibration.}
The relative cofibrancy of \(F_*\to\mathcal A_{\mathrm{fin}}\)
from \cite[Proposition~5.3.7]{Henry2016} therefore gives a section
of \(p\) sending the marked \(D_0\) to its cylinder \((D_0,D_1)\).
The lifting argument of \cite[Lemma~5.3.9]{Henry2016} then gives
the asserted cylinder-coherator property, as in
\cite[Lemma~5.3.11 and Corollary~5.3.12]{Henry2016}.

Henry's Theorem~5.1.2(5) now gives a fundamental-infinity-groupoid
right Quillen equivalence from spaces to this completion, with
geometric realization as its left adjoint.  By the second step the
completion is the original \(\Mod(\mathcal C)\) with its canonical
weak equivalences.  This proves the displayed statements and is the
implication summarized in \cite[Corollary~5.3.13]{Henry2016}.
\end{proof}
